% Part: first-order-logic
% Chapter: sequent-calculus
% Section: provability-consistency

\documentclass[../../../include/open-logic-section]{subfiles}

\begin{document}

\iftag{FOL}
      {\olfileid[id]{fol}{seq}{prv}}
      {\olfileid[id]{pl}{seq}{prv}}

\olsection{\usetoken{S}{derivability} dan Konsistensi}

Sekarang kita akan membuktikan sejumlah sifat relasi !!{derivability}.
Sifat-sifat tersebut menarik secara mandiri, tetapi masing-masing juga
akan berperan dalam pembuktian teorema kelengkapan.

\begin{prop}\ollabel{prop:provability-contr}
  Jika $\Gamma \Proves !A$ dan $\Gamma \cup \{!A\}$ tak konsisten,
  maka $\Gamma$ tak konsisten.
\end{prop}

\begin{proof}
Terdapat himpunan berhingga $\Gamma_0$ dan $\Gamma_1 \subseteq \Gamma$
sedemikian sehingga $\Log{LK}$ !!{derive} $\Gamma_0 \Sequent !A$ dan
$!A, \Gamma_1 \Sequent \quad$. Misalkan $\Log{LK}$-!!{derivation} untuk
$\Gamma_0 \Sequent !A$ adalah~$\pi_0$ dan $\Log{LK}$-!!{derivation}
untuk $\Gamma_1, !A \Sequent \quad$ adalah~$\pi_1$. Kita kemudian dapat
!!{derive}
\begin{prooftree}
\AxiomC{}
\RightLabel{$\pi_0$}
\Deduce$ \Gamma_0 \fCenter !A $
\AxiomC{}
\RightLabel{$\pi_1$}
\Deduce$!A, \Gamma_1 \fCenter $
\RightLabel{\Cut}
\BinaryInf$ \Gamma_0,\Gamma_1 \fCenter $
\end{prooftree}

Karena $\Gamma_0 \subseteq \Gamma$ dan $\Gamma_1 \subseteq \Gamma$,
maka $\Gamma_0 \cup \Gamma_1 \subseteq \Gamma$, sehingga $\Gamma$ tak
konsisten.
\end{proof}

\begin{prop}
\ollabel{prop:prov-incons}
$\Gamma \Proves !A$ jika dan hanya jika $\Gamma \cup \{\lnot !A\}$
tak konsisten.
\end{prop}

\begin{proof}
Pertama, andaikan $\Gamma \Proves !A$, yakni terdapat
!!a{derivation}~$\pi_0$ untuk $\Gamma \Sequent !A$. Dengan menambahkan
aturan $\LeftR{\lnot}$, kita memperoleh !!a{derivation} untuk
$\lnot !A, \Gamma \Sequent \quad$, yakni $\Gamma \cup \{\lnot !A\}$
tak konsisten.

Jika $\Gamma \cup \{\lnot !A\}$ tak konsisten, terdapat
!!a{derivation}~$\pi_1$ untuk $\lnot !A, \Gamma \Sequent \quad$.
Berikut ini adalah !!a{derivation} untuk $\Gamma \Sequent !A$:
\begin{prooftree}
  \Axiom$!A \fCenter !A$
  \RightLabel{\RightR{\lnot}}
  \UnaryInf$\fCenter !A, \lnot !A$
  \AxiomC{}
  \RightLabel{$\pi_1$}
  \Deduce$\lnot !A, \Gamma \fCenter$
  \RightLabel{\Cut}
  \BinaryInf$\Gamma \fCenter !A$
\end{prooftree}
\end{proof}

\begin{prob}
Buktikan bahwa $\Gamma \Proves \lnot !A$ jika dan hanya jika
$\Gamma \cup \{!A\}$ tak konsisten.
\end{prob}

\begin{prop}\ollabel{prop:explicit-inc}
  Jika $\Gamma \Proves !A$ dan $\lnot !A \in \Gamma$, maka $\Gamma$
  tak konsisten.
\end{prop}

\begin{proof}
  Andaikan $\Gamma \Proves !A$ dan $\lnot !A \in \Gamma$. Maka terdapat
  !!a{derivation}~$\pi$ untuk suatu sekuen $\Gamma_0 \Sequent !A$.
  Sekuen $\lnot !A, \Gamma_0 \Sequent \quad$ juga !!{derivable}:
  \begin{prooftree}
    \AxiomC{}
    \RightLabel{$\pi$}
    \Deduce$\Gamma_0 \fCenter !A$
    \Axiom$!A \fCenter !A$
    \RightLabel{\LeftR{\lnot}}
    \UnaryInf$\lnot !A, !A \fCenter$
    \RightLabel{\LeftR{\Exchange}}
    \UnaryInf$!A, \lnot !A \fCenter$
    \RightLabel{\Cut}
    \BinaryInf$\Gamma_0, \lnot !A \fCenter$
  \end{prooftree}
  Karena $\lnot !A \in \Gamma$ dan $\Gamma_0 \subseteq \Gamma$, hal ini
  menunjukkan bahwa $\Gamma$ tak konsisten.
\end{proof}

\begin{prop}\ollabel{prop:provability-exhaustive}
  Jika $\Gamma \cup \{!A\}$ dan $\Gamma \cup \{\lnot !A\}$ keduanya
  tak konsisten, maka $\Gamma$ tak konsisten.
\end{prop}

\begin{proof}
Terdapat himpunan berhingga $\Gamma_0 \subseteq \Gamma$ dan
$\Gamma_1 \subseteq \Gamma$, serta $\Log{LK}$-!!{derivation} $\pi_0$
dan $\pi_1$ masing-masing untuk $!A, \Gamma_0 \Sequent \quad$ dan
$\lnot !A, \Gamma_1 \Sequent \quad$. Kita kemudian dapat !!{derive}
\begin{prooftree}
\AxiomC{}
\RightLabel{$\pi_0$}
\Deduce$ !A, \Gamma_0 \fCenter $
\RightLabel{\RightR{\lnot}}
\UnaryInf$ \Gamma_0 \fCenter \lnot !A$
\AxiomC{}
\RightLabel{$\pi_1$}
\Deduce$\lnot !A, \Gamma_1 \fCenter  $
\RightLabel{\Cut}
\BinaryInf$ \Gamma_0, \Gamma_1 \fCenter $
\end{prooftree}
Karena $\Gamma_0 \subseteq \Gamma$ dan $\Gamma_1 \subseteq \Gamma$,
maka $\Gamma_0 \cup \Gamma_1 \subseteq \Gamma$. Jadi, $\Gamma$ tak
konsisten.
\end{proof}

\end{document}
